\section{What remains open}\label{sec:open}

\subsection{The corner of the trapezoid}
Theorem B covers the critical side $P_1P_2$ of the trapezoid of \cite[(1.5)]{BRS25} except the corner
$P_2$, where $q=q_{\rm opt}(p)$. At $P_2$ the scale-by-scale estimates behind Theorem \ref{thm:brs} sum
only with a loss: the exponent $r_*(p,p_\circ)$ reaches the Knapp line $r=\frac{n-1}{n+1}p'$ only in the
limit \eqref{eq:rstarlimit}, and the hypothesis of Theorem \ref{thm:brs} asks for $r<r_*$. The method would
need $\VBR(p,\frac{n-1}{n+1}p')$, the endpoint of the conjectured range of these inequalities, which is
established in \cite{BRS25} only at $p=p_{\rm ST}$. For $p=p_{\rm ST}$ the corner is reached in
\cite[Theorem 1.5]{BRS25}, through the $L^2$ restriction theorem; for $n=2$ the corner is also open
except at $\lambda=\frac16$ \cite[\S1.2]{BRS25}.

\subsection{The other sides}
For $0<\lambda<\frac{n-1}2$, sparse bounds on the top side $P_2P_3$ of the trapezoid, away from the
critical line, are open for $n\ge3$, and in the plane except at $\lambda=\frac16$. They are a
strengthening of off-diagonal $L^p\to L^r$ bounds for Bochner--Riesz means below the critical index,
and Theorem \ref{thm:BR} does not reach them. By duality the side $P_3P_4$ is equivalent to the side
$P_1P_2$, so Theorem B also gives it, except for the corner $P_3$.

\subsection{The range of the weights}
The exponent $p_1(w)=1+\frac{n-1}{n+1}(1-\frac1\sigma)$ of Theorem C is the one produced by Theorem
\ref{thm:fn} from the sparse bounds of Theorem B. Whether it is the sharp range in terms of the reverse
H\"older exponent is not known, as already noted in \cite[\S1.1]{BRS25}; weighted weak-type bounds
fail for $p>1$ for the power weights $|x|^\eps$, $\eps>0$ \cite[\S1.1]{BRS25}, so some restriction of
this kind is necessary.

\subsection{Other domains}
Parts I and II are written for the round sphere. Beltran, Roos and Seeger work with the Minkowski
functional $\rho$ of a convex domain with smooth boundary of positive curvature, and their
Theorem \ref{thm:brs} holds in that generality. The weak-type endpoint for such $\rho$ in dimensions
$n\ge3$ would follow in the same way from the Bochner--Riesz conjecture for $(1-\rho(\xi))^\lambda_+$,
which the local smoothing argument of Part I does not address.

\subsection{Maximal functions}
The maximal operator $\sup_{t>0}|R^{\lambda(p)}_{a,t}f|$ is not treated here. Already for
$\lambda>\lambda(p)$ and $p<2$ its $L^p$ boundedness requires conditions beyond those of the
Bochner--Riesz conjecture \cite{Tao02}, and nothing in this paper bears on it.
